% \documentclass[10pt, twocolumn, a4paper]{article}
\documentclass[12pt,a4paper]{article}

\usepackage[backend=biber, style=ieee]{biblatex}                        % To include the bibliography
\usepackage[left=2cm, right=2cm, top=2.5cm, bottom=2.5cm]{geometry}     % To set the margins
\usepackage[noend]{algpseudocode}
\usepackage[table]{xcolor}                                              % For coloring cells

\usepackage{algorithm}                                                  % To include algorithms
\usepackage{amsfonts}                                                   % To include math fonts:ToggleTerm direction=float
\usepackage{amsmath}                                                    % To include Mathematic symbols
\usepackage{authblk}                                                    % To format author affiliations
\usepackage{caption}                                                    % For caption spacing
\usepackage{enumitem}                                                   % To customize lists (items like i, ii, iii, iv)
\usepackage{float}                                                      % To place figures where you want them
\usepackage{fancyhdr}                                                   % To customize headers and footers
\usepackage{graphicx}                                                   % To include images
\usepackage{hyperref}                                                   % To include hyperlinks
\usepackage{lipsum}                                                     % TODO: remove this
\usepackage{listings}                                                   % To include code
\usepackage{tabularx}                                                   % For equal-width columns
\usepackage{tcolorbox}                                                  % To make colored boxes
\usepackage{tikz}                                                       % To draw graphs
\usepackage{titlesec}                                                   % To format section titles
\usepackage{xcolor}                                                     % To define colors

\usetikzlibrary{graphs,graphs.standard}
\usetikzlibrary{positioning}

\addbibresource{./references.bib}

% Esto es para poder hacer cajitas de código con el fondo gris
\lstset{
    language=C++,
    basicstyle=\ttfamily\footnotesize,
    keywordstyle=\color{blue}\bfseries,
    stringstyle=\color{green!60!black},
    commentstyle=\color{gray},
    backgroundcolor=\color{gray!05},
    frame=single,
    numbers=left,
    numberstyle=\footnotesize,
    stepnumber=1,
    numbersep=10pt,
    tabsize=2,
    showstringspaces=false,
    captionpos=b,
    breaklines=true,
}

% Para poder hacer flechas
\usetikzlibrary{shapes, arrows}

% Sección de definiciones
\titleformat{\section}{\Large\bfseries}{\thesection}{1em}{}
\titleformat{\subsection}{\large\bfseries}{\thesubsection}{1em}{}

% Caja de colores
\definecolor{mint}{RGB}{202,251,202}
\definecolor{yellow}{RGB}{255,255,202}
\definecolor{red}{RGB}{255,202,202}

% Variables globales para el documento
\newcommand{\facultyname}{UP | Facultad de Ingeniería}
\newcommand{\coursename}{Programación con GPUs}
\newcommand{\currentsemester}{Segundo Semestre}
\newcommand{\currentyear}{2025}

% Esto es para poder agregar comentarios al código
\newcommand{\comentario}[1]{\textcolor{gray}{// #1}}

% Definir los encabezados y pies de página
\pagestyle{fancy}
\fancyhf{} % Borra encabezados y pies de página

% Header configuration
\fancyhead[R]{\includegraphics[height=50px]{../apuntes/images/logo_up.jpg}}
\renewcommand{\headrulewidth}{0.4pt} % Agregar línea debajo del encabezado
\setlength{\headheight}{50pt} % Ajustar la altura del encabezado

% Footer configuration
\fancyfoot[L]{\facultyname}
\fancyfoot[C]{\thepage}
\fancyfoot[R]{\coursename}
\renewcommand{\footrulewidth}{0.4pt} % Agregar línea arriba del pie de página

% Configuración de la primera página
\fancypagestyle{firstpage}{
  \fancyhf{} % Borra encabezados y pies de página
  \fancyfoot[C]{\thepage} % Número de página centrado
  \renewcommand{\headrulewidth}{0pt} % Sin línea en el encabezado
  \renewcommand{\footrulewidth}{0.4pt} % Sin línea en el pie de página
}

% Configuración de la primera página
\AtBeginDocument{\thispagestyle{firstpage}}


\begin{document}

\begin{center}
  \LARGE\textbf{\coursename} \\
  \Large{Práctica 02 - Introducción a CUDA} \\
  \normalsize{\currentsemester, \currentyear} \\
  \vspace{1em}
  \hrule
\end{center}

\setcounter{section}{2}

\input{../shared.tex/soluciones_warning_box.tex}

\newpage

\tableofcontents

\newpage

\subsection{Ejercicio: Suma de dos vectores}

Si quisiéramos utilizar un thread para calcular la suma de los dos vectores. ¿Cómo se calcularía el índice del thread?

\begin{enumerate}
  \item \textbf{PISTA 1}: El índice del thread se calcula teniendo en cuenta el tamaño del bloque, el número de bloque y
    el número de thread.

  \item \textbf{PISTA 2}: El índice del thread se mueve entre 0 y el tamaño del bloque menos uno. El número de bloque
    se mueve entre 0 y el número de bloques menos uno.

\end{enumerate}

\textbf{SOLUCIÓN}: El índice del thread se calcula como:

\begin{equation}
  \texttt{idx = blockIdx.x * blockDim.x + threadIdx.x;}
\end{equation}

donde \texttt{blockIdx.x} es el número de bloque, \texttt{blockDim.x} es el tamaño del bloque y \texttt{threadIdx.x} es
el número de thread dentro del bloque.

\subsection{Ejercicio: Sumar elementos adyacentes}

Supongamos ahora que queremos usar un thread para sumar dos elementos contiguos de un vector (por simplicidad podemos
pensar que el vector y el de bloque tiene una cantidad par de elementos). ¿Cómo se calcularía el índice del thread?

\begin{enumerate}
  \item \textbf{PISTA 1}: El índice del thread tiene que ir de dos en dos. 

  \item \textbf{PISTA 2}: Usá una hoja y un lápiz / lapicera para calcular el índice del thread a mano a ver cuál podría
    ser.
\end{enumerate}

\textbf{SOLUCIÓN}: El índice del thread se calcula como: \texttt{idx = (blockIdx.x * blockDim.x + threadIdx.x) * 2;}.

\subsection{Ejercicio: Calcular el tamaño de los threads}

\begin{enumerate}
  \item \textbf{PISTA 1}: Primero hay que calcular cuántos bloques de threads serían necesarios.

  \item \textbf{PISTA 1}: Dado que 1024 no es divisor de 8000 siempre vamos a tener más threads que elementos a
    procesar. ¿Cómo se calcularía la cantidad de bloques sabiendo esa información?
\end{enumerate}

\textbf{SOLUCIÓN}: El total de threads es \textbf{8192}, que se calcula de la siguiente manera:

\begin{equation}
  numBlocks = \left\lceil \frac{8000}{1024} \right\rceil * 1024 = 8 * 1024 = 8192
\end{equation} \\

\subsection{Ejercicio: cudaMalloc (parte 1)}

\begin{enumerate}
  \item \textbf{PISTA 1}: El primer argumento de \texttt{cudaMalloc} es el puntero a la memoria que queremos reservar.

  \item \textbf{PISTA 2}: El segundo argumento de \texttt{cudaMalloc} es el tamaño de la memoria que queremos reservar.
\end{enumerate}

\textbf{SOLUCIÓN}: El segundo argumento de \texttt{cudaMalloc} es el tamaño de la memoria que queremos reservar, que se
calcula como \footnote{\href{https://docs.nvidia.com/cuda/cuda-runtime-api/group__CUDART__MEMORY.html}{CUDA
documentation}}:

\begin{equation}
  \texttt{sizeof(int) * v}
\end{equation}

\subsection{Ejercicio: cudaMalloc (parte 2)}

\begin{enumerate}
  \item \textbf{PISTA 1}: El primer argumento de \texttt{cudaMalloc} es el puntero a la memoria que queremos reservar.

  \item \textbf{PISTA 2}: El segundo argumento de \texttt{cudaMalloc} es el tamaño de la memoria que queremos reservar.
\end{enumerate}

\textbf{SOLUCIÓN}: El primer argumento de \texttt{cudaMalloc} es el puntero a un puntero a la memoria que queremos
reservar, ya que es la única forma en C para poder modificar el puntero (hacer una referencia).
\footnote{\href{https://docs.nvidia.com/cuda/cuda-runtime-api/group__CUDART__MEMORY.html}{CUDA Documentation}}. La
expresión sería entonces:

\begin{equation}
  \texttt{(void**)\&d\_A}
\end{equation}

\subsection{Ejercicio: Copia de memoria desde el \textit{host} al \textit{device}}

\begin{enumerate}
  \item \textbf{PISTA 1}: La función \texttt{cudaMemcpy} tiene cuatro argumentos (pensá qué argumentos necesitarías para
    poder copiar esa información)

  \item \textbf{PISTA 2}: Los 3 primeros argumentos son fáciles de descubrir, porque son los punteros de la memoria de
    origen, la memoria de destino y el tamaño de la memoria a copiar. Siendo que \texttt{cudaMemcpy} es una función que
    funciona en ambos sentidos ¿cuál sería el cuarto argumento?.
\end{enumerate}

\textbf{SOLUCIÓN}: La llamada a la API apropiada para esta copia de datos en CUDA es:

\begin{equation}
  \texttt{cudaMemcpy((void*)d\_A, (void*)h\_A, 3000, cudaMemcpyHostToDevice)}
\end{equation}

\subsection{Ejercicio: Manejo de errores}

\begin{enumerate}
  \item \textbf{PISTA 1}: Las funciones de manejo de memoria de cuda devuelven siempre un tipo de error, por convención
    es un tipo de datos especial que termina con el sufijo \texttt{\_t} ¿te acordás cuál es?
\end{enumerate}

\textbf{SOLUCIÓN}: El tipo de error que devuelven las funciones de manejo de memoria de CUDA es \texttt{cudaError\_t}.
Con lo cual la correcta definición de una variable (\texttt{err}) para recibir el error de las funciones de manejo de
memoria de CUDA sería:

\begin{equation}
  \texttt{cudaError\_t err;}
\end{equation}

\subsection{Ejercicio: funciones en CUDA}

\begin{enumerate}
  \item \textbf{PISTA 1}: Las funciones de CUDA tienen modificadores para el host como para el device. ¿Te acordás cómo
    se llaman?

  \item \textbf{PISTA 2}: ¿Hay realmente que declararlas dos veces?
\end{enumerate}

\textbf{SOLUCIÓN}: Las funciones de CUDA tienen modificadores para el host como para el device, y estas son
\texttt{\_\_host\_\_} y \texttt{\_\_device\_\_}. En realidad no es necesario declarar las funciones dos veces, ya que el
compilador se encarga de generar el código para ambos dispositivos. Por ejemplo:

\begin{equation}
  \texttt{\_\_host\_\_ \_\_device\_\_ void foo(int *a) \{ ... \}}
\end{equation}

\input{../shared.tex/common_footers.tex}

\end{document}
